\begin{abstract}% .
% Pvsmthosofn onuprvisdf-u(SFT),whchco fuda-
% mentally enhhe model's capabilityddress this
% limittioby introdung a GRPO-bsed optimizationdireim-
% provmodl' reasnigabiity.Ou proacleves-
% toguide blingmore effective
% and ph-
% that ou methd achies superior-
% advtagGRPO over trdiionl SFT ppah
% 
% Reasoning over knowledge graphs is a fundamental challenge in artificial intel-
% ligence, requiring both accurate inference and interpretability. Recent advances
% leverage reinforcement learning to optimize reasoning paths, yet most approaches
% focus on end-state rewards, neglecting the importance of intermediate reasoning
% steps. In this work, we propose a deliberative and constructive reasoning frame-
% work that explicitly rewards the process of path exploration on knowledge graphs.
% By integrating process-based rewards into policy optimization, our method en-
% courages agents to discover interpretable and effective reasoning trajectories. Ex-
% periments on benchmark datasets demonstrate improved performance in both an-
% swer accuracy and path quality, highlighting the significance of path-aware policy
% design for knowledge graph reasoning.
% Large language models (LLMs) have demonstrated great potential in various reasoning tasks. 
% However, their performance in knowledge graph reasoning (KGR) remains suboptimal due to the lack of interpretability and the challenge of capturing complex relationships.
% In this work, we propose a novel framework that leverages process-based rewards to enhance the reasoning capabilities of LLMs on knowledge graphs. 
% By integrating these rewards into a reinforcement learning paradigm, 
% our approach encourages the model to explore more effective and interpretable reasoning paths.
%  Experimental results on benchmark datasets demonstrate that our method significantly improves both answer accuracy and path quality,
%   showcasing the advantages of process-aware optimization in KGR tasks.
% Knowledge Graphs(KGs) have demonstrated great potential in complex facilitating reasoning tasks due to their ability to capture intricate relationships
% and provide interpretable reasoning paths.
% Recent advances in the utilization of KG for large language models(LLMs) mainly focus on , igno
% Reasoning over knowledge graphs (KGs) is a critical task to
% improve the factuality and mitigate the hallucinations of large language models (LLMs).
% % Recent advances predominantly emphasize answer generation by decoding from KGs, often overlooking the necessity for deep understanding and robust reasoning capabilities of LLMs over knowledge graphs. 
% % Recent advances mainly focus on  
% % designing decoding rules to extract rule-compliant answers from knowledge graphs,
% % while neglecting deep understanding of graph features. 
% % This results in insufficient fundamental reasoning capabilities of large language models over knowledge graphs.
% % Recent advances mainly focus on 
% %  rule-based decoding for extracting answers from KGs,
% % while neglecting 
% % Consequently, the fundamental reasoning capacity 
% % of large language models over knowledge graphs remains underdeveloped. 
% Recent advances primarily decouple the generation of reasoning paths
% from the subsequent answer inference based on those paths,
% preventing LLMs from comprehensively understanding the original graph structure 
% and limiting their capacity for self-directed exploration.
% % In this paper, we proporse EOG (Exploration on Graphs), a novel framework that 
% In this paper, we propose a novel end-to-end framework that encourages LLMs to actively
% explore on the graphs (EOG) to enable spontaneous and deliberative reasoning.
% First, we model the graph reasoning process as a textual chain of thought to 
% utilize the reasoning capabilities of LLMs. 
% Then, we introduce the Group Relative Policy Optimization (GRPO) to 
% enable the LLMs to learn the experience of the spontaneous exploration on the graphs.
% % from a group of reasoning paths and 
% % and refine their exploration policy towards the correct answer.
% % to produce diverse reasoning thoughts, incentivizing 
% % the autonomous exploration of the KGs.
% % Moreover, we proporse a path-based rewards aligned with 
% % To generate more effective and faithful reasoning paths, we propose 
% % the path-based rewards based on the quality of the reasoning paths.
% Extensive experiments on three benchmark datasets demonstrate that our method
% significantly outperforms state-of-the-art baselines with different open-source LLMs,
% showing the promising advantages of path-wise reward modeling in KG reasoning tasks.
% Reasoning over knowledge graphs (KGs) is a critical task to
% improve the factuality and mitigate the hallucinations of large language models (LLMs).
% Recent advances predominantly emphasize answer generation by decoding from KGs, often overlooking the necessity for deep understanding and robust reasoning capabilities of LLMs over knowledge graphs. 
% Recent advances mainly focus on  
% designing decoding rules to extract rule-compliant answers from knowledge graphs,
% while neglecting deep understanding of graph features. 
% This results in insufficient fundamental reasoning capabilities of large language models over knowledge graphs.
% Recent advances mainly focus on 
%  rule-based decoding for extracting answers from KGs,
% while neglecting 
% Consequently, the fundamental reasoning capacity 
% of large language models over knowledge graphs remains underdeveloped. 
% Knowledge Graph Question Answering (KGQA) is a critical task to
% improve the factuality and mitigate the hallucinations of large language models (LLMs).\
% Enabling Large Language Models (LLMs) to reason over knowledge graphs 
% is a crucial approach for Knowledge Graph Question Answering (KGQA), 
% aiming to improve their factuality and mitigate hallucinations.
% Reasoning over knowledge graphs is a crucial approach for Large Language Models (LLMs) 
% to improve their factuality and mitigate hallucinations 
% in Knowledge Graph Question Answering (KGQA).
% Knowledge Graph Question Answering (KGQA) is a critical task to
% improve the factuality and mitigate the hallucinations of large language models (LLMs),
% which involves complex reasoning over structured knowledge graphs.
% Reasoning over knowledge graphs is a crucial approach for Large Language Models (LLMs) to 
% improve their factuality and mitigate hallucinations.
% Although Large language models (LLMs) have shown remarkable capabilities,
The reasoning process of  Large Language Models (LLMs) is often plagued by hallucinations and missing facts in question-answering tasks.
A promising solution is to ground LLMs' answers in verifiable knowledge sources, such as Knowledge Graphs (KGs).
% However, recent KG-enhanced methods 
% primarily decouple the generation of reasoning paths
% from the subsequent answer inference based on those paths,
% % preventing LLMs from comprehensively understanding the original graph structure 
% % and limiting their capacity for active exploration of the knowledge graph.
% preventing LLMs from holistically leveraging and exploring the original graph structure to derive the final answer.
% In this paper, we proporse EOG (Exploration on Graphs), a novel framework that 
% In this paper, we propose a novel synergistic exploitation-exploration framework 
% to enable structured and spontaneous reasoning on knowledge graphs (SSR-KG).
% (i) In the exploitation stage, 
% we model the graph reasoning process as the generation of textual chain of thought 
% to enhance the structure and completeness of the reasoning process through the supervised fine-tuning (SFT).
% % and 
% % and utilize the reasoning capabilities of LLMs by supervised fine-tuning (SFT).
% (ii) In the exploration stage, we propose the adaptive path-wise reward to score the quality
% of the reasoning paths, thereby dynamically refining the LLMs' exploration policy 
% to promote more spontaneous and efficient exploration by the reinforcement learning (RL).
% from a group of reasoning paths and 
% and refine their exploration policy towards the correct answer.
% to produce diverse reasoning thoughts, incentivizing 
% the autonomous exploration of the KGs.
% Moreover, we proporse a path-based rewards aligned with 
% To generate more effective and faithful reasoning paths, we propose 
% the path-based rewards based on the quality of the reasoning paths.
% 目前的方法主要通过预定的 学习过的样本
% Recent KG-enhanced methods mainly focus on constructing structured schemas and instructions to govern LLMs' reasoning process.
% However, their rigid and pre-defined nature limits the model's flexibility and exploration capabilities in complex reasoning scenarios.
% Recent KG-enhanced methods mainly focus on constructing structured schemas or learning fixed features from known samples to govern LLMs' reasoning process.
Prevailing KG-enhanced methods typically constrained LLM reasoning either by enforcing rules during generation or by imitating paths from a fixed set of demonstrations.
% However, these approaches limits the model's flexibility and exploration capabilities, leading to insufficient utilization of graphs and thus impairing the model's generalization in complex reasoning scenarios.
% However, these approaches restricted LLMs to previously seen patterns, fundamentally limiting their generalization to complex graphs with unseen reasoning structures.
% this reliance on static and pre-defined patterns fundamentally limit the model's ability to explore unseen paths of the graph, thus impairing its generalization to more complex and novel reasoning paths.
% However, these approaches restrict LLMs to previously seen patterns, fundamentally limiting their generalization to complex graphs. 
However, they naturally confined the reasoning patterns of LLMs within the scope of prior experience or fine-tuning data, limiting their generalizability to out-of-distribution graph reasoning problems.
% To tackle this problem,
To address this issue,
% , we aim to enhance LLMs' reasoning capability by encouraging to explore a more diverse reasoning space.
% To overcome this, our goal is exposing the model to a broader and richer reasoning patterns
% is to promote the generalization through autonomous exploration, exposing the model to a broader and more complex set of reasoning patterns.
% they usually suffered from the inherent inflexibility of these prescribed paradigms,
% hindering their ability to actively explore the graph for alternative and potentially more effective reasoning process.
% restricts the model to the demonstrated patterns and discourages it from discovering alternative, potentially more efficient or novel, reasoning strategies.
in this paper, we propose Explore-on-Graph (EoG), a novel framework that encourages LLMs to autonomously explore a more diverse reasoning space on KGs.
To incentivize exploration and discovery of novel reasoning paths, we propose to introduce reinforcement learning during training, whose reward is the correctness of the reasoning paths' final answers. 
To enhance the efficiency and meaningfulness of the exploration, we propose to incorporate path information as additional reward signals to refine the exploration process and reduce futile efforts.
Extensive experiments on five KGQA benchmark datasets demonstrate that, to the best of our knowledge, our method achieves state-of-the-art performance, outperforming not only open-source but also even closed-source LLMs\footnote{Code and data are available at: https://github.com/ysq111333/EoG}. 
% Our code is publicly available at xx.
% promotes the autonomous exploration of LLMs on knowledge graphs through path-refined reward modeling.
% To effectively optimize the LLM’s exploration policy, we introduce the Group Relative Policy Optimization (GRPO), which leverages the group-wise relative advantages to facilitate stable and efficient policy updates. 
% To provide fine-grained feedback and mitigate model hallucinations, we propose path-refined rewards that effectively guide the LLM’s exploration process towards factually grounded paths.
% EoG introduces path-refined rewards that effectively guide the LLM’s exploration process towards factually grounded paths.
% providing fine-grained feedback and mitigate model hallucinations.
% Specifically, EoG adopts an answer-oriented reward to incentivize LLMs to explore diverse reasoning paths that lead to potentially correct answers and a path-refinement reward which evaluates the logical coherence to reduce hallucinations and enhance the overall quality of the reasoning process.
% The following advantages of EoG are examined and illustrated through a number of well-designed experiments: 
% 1) EOG achieves deep exploration on knowledge graphs by LLMs.
% 2) EOG demonstrates outstanding performance in terms of reasoning depth and comprehensiveness.
% 3) EOG exhibits robust generalization capabilities across diverse datasets.
% First, we model graph reasoning as the generation of long chain-of-thought to leverage the reasoning capabilities of LLMs. 
% we introduce the Group Relative Policy Optimization (GRPO) to incentivize LLMs to learn from a diverse range of explorations.
% % and to autonomously distill effective reasoning strategies from these experiences.
% % enable the LLMs to produce diverse reasoning thoughts and , incentivizing the autonomous exploration on the KGs.
% Moreover, 
% we propose path-refined rewards to effectively guide the model's exploration
% % by providing a more fine-grained and reasonable exploration experiences 
% based on the entire reasoning paths.
% Extensive experiments on five KGQA benchmark datasets
% demonstrate that our method achieves state-of-the-art performance, outperforming even closed-source large-scale models.
\end{abstract}
